\section*{Chapitre \ref{ch:b3:ocean-dynamics} --- Dynamique océanique}

\begin{solution}{exo:b3:ocean-dynamics:1}
$f = \qty{9.37e-5}{s^{-1}}$; $\tau/(\rho f) = 0.12/(1025\times\num{9.37e-5})
= \qty{1.25}{m^2/s}$, à droite du vent~: vers le sud. Sur
\qty{1000}{km}, \qty{1.25e6}{m^3/s}, \qty{1.25}{\sverdrup}.
\end{solution}

\begin{solution}{exo:b3:ocean-dynamics:2}
$f = \qty{1.031e-4}{s^{-1}}$; $D_E = \pi\sqrt{2\times0.05/\num{1.031e-4}}
= \qty{98}{m}$. Le courant y vaut $\mathrm e^{-\pi} = \num{0.043}$ du
courant de surface, circulant en sens inverse.
\end{solution}

\begin{solution}{exo:b3:ocean-dynamics:3}
$\beta = 2\Omega\cos\ang{30}/a = \qty{1.98e-11}{m^{-1}.s^{-1}}$; $V =
\num{-1e-7}/(1025\times\num{1.98e-11}) = \qty{-4.92}{m^2/s}$~; sur
\qty{5000}{km}, \qty{-24.6}{\sverdrup}~: vers le sud.
\end{solution}

\begin{solution}{exo:b3:ocean-dynamics:4}
Les deux \omterm{def:b3:ocean-dynamics:tides}{composantes de marée} se retrouvent en phase après $1/(1/12 - 1/12.4206) =
\qty{354}{h}$, \num{14.77}\text{ jours}, soit un demi-mois synodique.
\end{solution}

\begin{solution}{exo:b3:ocean-dynamics:5}
À \ang{40} N, $f = \qty{9.37e-5}{s^{-1}}$, $\beta =
\qty{1.75e-11}{m^{-1}.s^{-1}}$, $\tau_x = \qty{0.05}{Pa}$ et
$\partial\tau_x/\partial y = 0.1/\num{1.11e6} = \qty{9.0e-8}{Pa/m}$. Le
premier terme, $-\num{9.0e-8}/(1025\times\num{9.37e-5}) =
\qty{-9.38e-7}{m/s}$~; le second terme, $0.05\times\num{1.75e-11}/(1025\times
(\num{9.37e-5})^2) = \qty{+9.7e-8}{m/s}$. $w_E = \qty{-8.4e-7}{m/s}$,
\qty{-26.5}{m/an}~: l'eau est poussée vers le bas.
\end{solution}

\begin{solution}{exo:b3:ocean-dynamics:6}
$\delta_S = \num{2e-6}/\num{2e-11} = \qty{100}{km}$~; vitesse $=
\num{30e6}/(\num{1e5}\times800) = \qty{0.375}{m/s}$.
\end{solution}

\begin{solution}{exo:b3:ocean-dynamics:7}
$\beta = \qty{1.98e-11}{m^{-1}.s^{-1}}$; $\delta_M = (10^4/\num{1.98e-11})^{1/3}
= \qty{79.6}{km}$, et le jet mesure $3.628\,\delta_M = \qty{289}{km}$ de large.
Avec $A_H$ huit fois plus faible, les deux longueurs sont divisées par deux, puisque $8^{1/3} = 2$.
\end{solution}

\begin{solution}{exo:b3:ocean-dynamics:8}
$f = \qty{8.57e-5}{s^{-1}}$; $v = (21.3 - 14.5)/(\num{8.57e-5}\times\num{8e4})
= \qty{0.99}{m/s}$, positif, $B$ étant à droite~: en regardant de $A$ vers $B$,
vers l'est, le courant se dirige vers le nord. La surface de la mer se trouve
$6.8/9.81 = \qty{0.69}{m}$ plus haute en $B$~; la colonne d'eau en $B$ est la
plus légère et plus chaude, à droite du courant.
\end{solution}

\begin{solution}{exo:b3:ocean-dynamics:9}
Marnage de vive-eau $1.5K(1 + 0.46) = \qty{0.78}{m}$~; marnage de morte-eau $1.5K(1 - 0.46)
= \qty{0.29}{m}$~; rapport $1.46/0.54 = 2.7$.
\end{solution}

\begin{solution}{exo:b3:ocean-dynamics:10}
Comme au \cref{thm:b3:ocean-dynamics:kelvin}, $u = \pm(g/c)\eta$ pour $F(x \mp
ct)$, et $fu = -g\,\partial\eta/\partial y$ devient $\beta_0y(\pm g/c)G =
-gG'$, donc $G = \exp(\mp\beta_0y^2/2c)$. Seul le signe supérieur, celui de l'onde
vers l'est, est borné des deux côtés de l'équateur. $\beta_0 =
\qty{2.29e-11}{m^{-1}.s^{-1}}$, $\sqrt{2.5/\num{2.29e-11}} = \qty{330}{km}$.
\end{solution}

\begin{solution}{exo:b3:ocean-dynamics:11}
(a)~En intégrant $\beta v = f\,\partial w/\partial z$ depuis le fond, $w = 0$,
jusqu'au sommet de la couche abyssale, $w = W_0$~: $\beta V = fW_0$, et $f/\beta =
a\tan\varphi$. (b)~$V = \num{6.371e6}\times0.577\times10^{-7} =
\qty{0.368}{m^2/s}$, \qty{1.84}{\sverdrup} sur \qty{5000}{km}, vers le pôle.
(c)~Issu d'un courant profond de bord ouest, alimenté par les sources de hautes latitudes~;
l'intérieur s'écoule vers les sources, et le courant de bord ouest s'en
éloigne.
\end{solution}

\begin{solution}{exo:b3:ocean-dynamics:12}
(a)~Avec $z$ orienté vers le haut à partir du fond, $A_zW'' = \mathrm if(W - W_g)$, le
gradient de pression étant celui de l'écoulement géostrophique au-dessus. $W - W_g =
C\mathrm e^{-(1+\mathrm i)z/\delta}$ est borné vers le haut, et $W(0) = 0$ donne
$C = -W_g$. (b)~Le défaut de transport vaut $\int_0^\infty(W - W_g)\,\mathrm dz
= -W_g\delta/(1+\mathrm i) = -\tfrac12\delta(1 - \mathrm i)W_g$, soit
$-\tfrac12\delta(u_g + v_g,\ v_g - u_g)$. L'écoulement géostrophique étant
non divergent, $w = -\nabla\cdot(\text{défaut}) = \tfrac12\delta(\partial_x
v_g - \partial_y u_g) = \tfrac12\delta\zeta_g$. (c)~$\delta =
\sqrt{0.02/10^{-4}} = \qty{14.1}{m}$; $r = 10^{-4}\times14.1/8000 =
\qty{1.77e-7}{s^{-1}}$, un temps d'amortissement de \num{65}\text{ jours}.
\end{solution}

\begin{solution}{pb:b3:ocean-dynamics:1}
\textbf{1.} $\tau_x = \qty{-0.1}{Pa}$ (d'est), 0 et \qty{+0.1}{Pa}
(d'ouest).

\textbf{2.} $f = \num{4.99e-5}$, \num{8.37e-5} et
\qty{1.117e-4}{s^{-1}}.

\textbf{3.} $V_E = -\tau_x/(\rho f)$~: \qty{+1.96}{m^2/s} (nord) à
\ang{20} N et \qty{-0.87}{m^2/s} (sud) à \ang{50} N. Ils convergent.

\textbf{4.} $D_E = \pi\sqrt{2\times0.05/\num{8.37e-5}} = \qty{109}{m}$.

\textbf{5.} $\tau_x = 0$ à \ang{35} N, donc $w_E = -(1/\rho f)\,
\partial\tau_x/\partial y = -\pi\tau_0/(\rho fL) = \qty{-1.10e-6}{m/s}$,
\qty{-34.7}{m/an}, vers le bas.

\textbf{6.} $500/34.7 = \num{14.4}\text{ ans}$.

\textbf{7.} $L = a\times\pi/6 = \qty{3336}{km}$~; $\beta =
\qty{1.875e-11}{m^{-1}.s^{-1}}$.

\textbf{8.} $\operatorname{curl}_z\vec\tau = -\pi\tau_0/L =
\qty{-9.42e-8}{Pa/m}$.

\textbf{9.} $V = \num{-9.42e-8}/(1025\times\num{1.875e-11}) =
\qty{-4.90}{m^2/s}$, vers le sud.

\textbf{10.} $4.90\times\num{6e6} = \qty{2.94e7}{m^3/s}$,
\qty{29.4}{\sverdrup} vers le sud.

\textbf{11.} $\qty{4.9}{mm/s}$~; \qty{1000}{km} en \qty{2.0e8}{s},
\num{6.5}~ans.

\textbf{12.} $fw_E/\beta = \num{8.37e-5}\times(\num{-1.10e-6})/
\num{1.875e-11} = \qty{-4.90}{m^2/s}$~: c'est la même valeur, car le
transport d'Ekman est nul là où $\tau_x = 0$. Une colonne aplatie au taux $w_E$
conserve sa \omterm{def:b3:atmospheric-dynamics:pv}{vorticité potentielle} $f/h$ seulement en se déplaçant vers un $f$ plus petit, c'est-à-dire vers le sud,
au taux $\beta V_g = fw_E$.

\textbf{13.} Une couche limite ne peut être qu'une couche dans laquelle la correction à
l'intérieur décroît en s'éloignant de la paroi~; avec $\beta > 0$ cela correspond à
$\mathrm e^{-x/\delta_S}$, qui décroît vers l'est, et seulement à partir de la paroi ouest.
En termes de vorticité, le frottement dans un courant vers le nord contre une
paroi est retirerait de la vorticité positive, ce qui est le mauvais signe.

\textbf{14.} $\delta_S = \num{e-6}/\num{1.875e-11} = \qty{53.3}{km}$.

\textbf{15.} $A = \qty{4.90}{m^2/s}$, $X/\delta_S = 112.5$:
$4.90\times111.5 = \qty{546}{m^2/s}$, \qty{0.55}{m/s} à travers
\qty{1000}{m}.

\textbf{16.} Là où $(X/\delta_S)\mathrm e^{-x/\delta_S} = 1$~: $x =
\delta_S\ln(X/\delta_S) = \qty{252}{km}$.

\textbf{17.} $A(X - \delta_S - x_0) = 4.90\times(6000 - 53 - 252)\times
10^3 = \qty{27.9}{\sverdrup}$. L'écoulement intérieur va encore vers le sud dans
la bande voisine de la paroi~; le déficit, $A(x_0 + \delta_S) =
\qty{1.5}{\sverdrup}$, est le transport intérieur qui aurait traversé
cette bande.

\textbf{18.} $\delta_M = (\num{5e3}/\num{1.875e-11})^{1/3} = \qty{64.4}{km}$;
jet $3.628\times64.4 = \qty{233}{km}$.

\textbf{19.} $1.163\times29.4 - A\times\qty{233}{km} = 34.2 - 1.1 =
\qty{33.0}{\sverdrup}$~; l'excédent, \qty{16}{\percent} du transport
intérieur, retourne vers le sud dans le contre-courant juste à l'est du jet.

\textbf{20.} $c = \sqrt{9.81\times4000} = \qty{198}{m/s}$; $L_R =
198/\num{8.37e-5} = \qty{2370}{km}$.

\textbf{21.} $198\times12.4206\times3600 = \qty{8860}{km}$.

\textbf{22.} Périmètre $2(6000 + 3336) = \qty{18670}{km}$, \qty{26.2}{h},
2.1 M$_2$ périodes.

\textbf{23.} Le cercle a pour longueur $2\pi a\cos\ang{35} = \qty{32790}{km}$, parcouru
en un jour lunaire, \qty{24.84}{h}~: \qty{367}{m/s}, ce qui demande $c^2/g =
\qty{13.7}{km}$ d'eau.

\textbf{24.} Avec un \omterm{def:b3:gfd-foundations:rossby-radius}{rayon de déformation de Rossby} comparable à la largeur du bassin
et la longueur d'onde de M$_2$ comparable à sa longueur, environ un
\omterm{def:b3:ocean-dynamics:amphidrome}{point amphidromique}, près du centre du bassin~; la pleine mer tourne dans le sens antihoraire
autour de lui, en laissant la côte à sa droite.

\textbf{25.} $\tfrac34\times0.357\times\cos^2\ang{35} = \qty{0.18}{m}$.

\textbf{26.} Le courant de bord ouest ramène vers le nord les
\textbf{\qty[text-series-to-math=true]{29}{\sverdrup}} que l'intérieur
transporte vers le sud (\qty{29.4}{\sverdrup} par le calcul).
\end{solution}
